\section{Conclusion}

Authority-aligned linearization rests on an operational fact that consensus and CRDT designs cannot assume: at every moment, one principal is designated as authoritative.
The incident-journal linearization point is the command edge, the vehicle edge node that currently holds that authority.
Reference data, field submissions, materialized state, and audit forwarding sit outside that boundary on weaker contracts.
Command transfer uses fenced promotion rather than leader election.

The model checks and the \systemname{} service-path prototype agree on the safety boundary.
Under strict fencing, no two writers and no forked journal; the weak variant produces the expected split-brain witness.
Under a responder partition, the command edge committed 20/20 writes.
The CRDT-LWW comparison shows the trade-off: CRDT replicas accept local writes at multiple nodes and converge quickly after anti-entropy exchange, while AAL preserves one incident journal and treats semantic conflict rejection as part of the linearization boundary.
The cost is no automatic failover and lower write availability for authority-bearing operations.

Deploy \systemname{} on a vehicle-mounted mesh, exercise fenced promotion across multiple command edges, implement service-path semantic rejection for authoritative conflicts, and add crash-boundary and Android end-to-end tests.
The harder comparison is empirical, not architectural: measure the design against CoNICE-class consensus and CRDT-based emergency systems on workloads where command authority is contested rather than designated.
